%
% 第8章：经典网络架构

\clearpage
\cleardoublepage

\chapter{经典网络架构}

\section{神经网络演进}

深度学习的历史显示了从简单网络到越来越复杂架构的发展，每个都在前一个时代的基础上进行创新。

\section{LeNet（1998）}

Yann LeCun的LeNet是首个成功用于数字识别的CNN。

\begin{architecture}[LeNet-5架构]
\begin{enumerate}
    \item 输入：$32\times32$ 灰度图像
    \item C1：$6$ 个 $5\times5$ 卷积核 $\rightarrow$ $28\times28\times6$
    \item S2：$2\times2$ 平均池化 $\rightarrow$ $14\times14\times6$
    \item C3：$16$ 个 $5\times5$ 卷积核 $\rightarrow$ $10\times10\times16$
    \item S4：$2\times2$ 平均池化 $\rightarrow$ $5\times5\times16$
    \item C5：$120$ 个 $5\times5$ 卷积核 $\rightarrow$ $120$
    \item F6：全连接，84个神经元
    \item 输出：10类（数字）
\end{enumerate}
\end{architecture}

参数：约60,000

关键创新：
\begin{itemize}
    \item 通过卷积实现权重共享
    \item 局部感受野
    \item 空间子采样（池化）
\end{itemize}

\section{AlexNet（2012）}

AlexNet以显著优势赢得ImageNet LSVRC-2012，彻底改变了计算机视觉。

参数：约6000万

关键创新：
\begin{itemize}
    \item ReLU激活（训练更快）
    \item GPU训练
    \item 数据增强
    \item Dropout正则化
    \item 响应归一化
\end{itemize}

\section{VGGNet（2014）}

VGGNet强调深度和简洁，仅使用 $3\times3$ 卷积。

参数：约1.38亿

关键洞察：两个 $3\times3$ 卷积的参数量比一个 $5\times5$ 卷积少约28\%。

\section{GoogLeNet / Inception（2014）}

引入了并行多尺度卷积的Inception模块。

关键创新：
\begin{itemize}
    \item 多尺度特征提取
    \item $1\times1$ 卷积降维
    \item 全局平均池化
    \item 辅助分类器
\end{itemize}

参数：约400万（非常高效！）

\section{ResNet（2015）}

ResNet引入了残差连接，实现了极深网络的训练。

关键创新：
\begin{itemize}
    \item 跳跃连接实现梯度流动
    \item 恒等捷径
    \item 瓶颈设计（$1\times1$, $3\times3$, $1\times1$）
    \item 可实现1000+层
\end{itemize}

参数：约2500万

\section{Transformer（2017）}

Transformer用自注意力替代了循环。

关键创新：
\begin{itemize}
    \item 自注意力替代循环
    \item 可并行化（不同于RNN）
    \item 捕获长期依赖
    \item 可随计算扩展
\end{itemize}

参数：65M（基础）到175B（GPT-3）及更多

\begin{table}[h]
    \centering
    \caption{架构比较}
    \begin{tabular}{L{2.5cm} L{2cm} L{2.5cm} L{2.5cm}}
        \toprule
        \textbf{网络} & \textbf{年份} & \textbf{参数} & \textbf{Top-1准确率} \\
        \midrule
        AlexNet & 2012 & 60M & 62.5\% \\
        VGG-16 & 2014 & 138M & 71.5\% \\
        GoogLeNet & 2014 & 4M & 72.0\% \\
        ResNet-50 & 2015 & 25M & 76.1\% \\
        ViT-B/16 & 2021 & 86M & 79.9\% \\
        \bottomrule
    \end{tabular}
\end{table}

\section{本章小结}

\begin{enumerate}
    \item LeNet开创了用于数字识别的CNN
    \item AlexNet通过GPU训练普及了深度学习
    \item VGG展示了通过 $3\times3$ 卷积实现深度
    \item Inception引入了多尺度处理
    \item ResNet通过跳跃连接实现了超深网络
    \item Transformer通过注意力革新了NLP并影响视觉领域
\end{enumerate}
